A quasi-plurisubharmonic function with zero Lelong numbers can still be
unbounded, but adding it to another potential changes no multiplier ideal.

\begin{lemma}\label{lem:zero-lelong-ideal-neutrality}\index{Multiplier ideal!zero Lelong residual}
Let $q$ be quasi-plurisubharmonic on a complex manifold and suppose that
$\nu(q,x)=0$ at every point. Then $\mathrm e^{-tq}$ is locally integrable
for every $t>0$. For every quasi-plurisubharmonic $v$ and every $t>0$,
\[
 \mathcal I\bigl(t(v+q)\bigr)=\mathcal I(tv).
\]
\end{lemma}
\begin{proof}
At each point $x$, Skoda's theorem \cite[Lemma~5.6]{DemAM11}, applied
after a local smooth correction, makes $\mathrm e^{-tq}$ integrable near
$x$, since $\nu(tq,x)=0<1$. Thus $q\sim_{\mathcal I}0$;
the addition rule in \cref{lem:hermitian-i-calculus}, whose proof is local, gives
$v+q\sim_{\mathcal I}v$, which is precisely the asserted equality at every weight.
\end{proof}
